\section{Results}
\label{sec:results}

\subsection{RQ1: accuracy against baselines}
% generated from results/e1_baselines/results.csv (sha256 of LF text 898355fc90852cc6, 300 rows)
% generated from results/e1b_scored_baselines/results.csv (sha256 of LF text 0d8b65d8aef12277, 300 rows)
\begin{table*}[t]
\centering
\caption{Ring-level detection on 5 scenarios $\times$ 10 test seeds, pooled (mean $\pm$ 95\% CI over runs;
a run with no flags has undefined precision and counts as F1 $=0$). TTD: median time from ring activation to
first flag. Latency: median per-batch detection time on the in-memory store; $^\dagger$full-graph
methods run on a 5\,s timer, not per batch, so their entry is the median time of one full-graph pass. $p$: paired Wilcoxon signed-rank
test of F1 against the proposed method on identical streams. Scored baselines use the shipped scorer (\texttt{models/ring\_scorer\_lr.json}), with the alert threshold selected per method on the validation seeds; the proposed method includes the same scorer, so only rows within a group isolate the candidate generator. $p$ values are Holm-corrected.}
\label{tab:baselines}
\begin{tabular}{lcccccc}
\toprule
Method & Precision & Recall & F1 & TTD (s) & Latency (ms) & $p$ \\
\midrule
\multicolumn{7}{l}{\textit{Without a scorer}} \\
Static Cypher (cap 100) & 0.12 $\pm$ 0.03 & 0.61 $\pm$ 0.08 & 0.19 $\pm$ 0.05 & 9.5 & 6.0$^\dagger$ & $<$0.001 \\
Static Cypher (cap 10) & 0.14 $\pm$ 0.04 & 0.61 $\pm$ 0.08 & 0.21 $\pm$ 0.05 & 9.5 & 5.9$^\dagger$ & $<$0.001 \\
Louvain communities & 0.39 $\pm$ 0.09 & 0.33 $\pm$ 0.08 & 0.29 $\pm$ 0.06 & 18.8 & 38.4$^\dagger$ & $<$0.001 \\
Connected components & -- & 0.00 $\pm$ 0.00 & 0.00 $\pm$ 0.00 & -- & 11.1$^\dagger$ & $<$0.001 \\
\midrule
\multicolumn{7}{l}{\textit{With the calibrated scorer}} \\
Static Cypher (cap 100) + scorer & 0.63 $\pm$ 0.08 & 0.48 $\pm$ 0.09 & 0.52 $\pm$ 0.08 & 8.6 & 4.5$^\dagger$ & $<$0.001 \\
Static Cypher (cap 10) + scorer & 0.64 $\pm$ 0.07 & 0.53 $\pm$ 0.09 & 0.56 $\pm$ 0.07 & 9.3 & 4.6$^\dagger$ & $<$0.001 \\
Louvain communities + scorer & 0.81 $\pm$ 0.12 & 0.32 $\pm$ 0.09 & 0.41 $\pm$ 0.09 & 19.0 & 27.9$^\dagger$ & $<$0.001 \\
Hub-aware incremental (proposed) & 0.76 $\pm$ 0.06 & 0.74 $\pm$ 0.06 & \textbf{0.73 $\pm$ 0.05} & 9.3 & 26.6 & -- \\
Proposed, a-priori thresholds & 0.73 $\pm$ 0.06 & 0.69 $\pm$ 0.06 & 0.68 $\pm$ 0.05 & 9.0 & 22.7 & 0.034 \\
\bottomrule
\end{tabular}
\end{table*}

\input{tables/tab1b_scenarios.tex}

Table~\ref{tab:baselines} pools \EScenarios{} scenarios $\times$ \ESeeds{}
test seeds for every method, all on identical streams. Without a scorer, the
static rules reach ring F1 \ECapTenF{} with a degree cap of 10 and
\ECapHundredF{} with a cap of 100, Louvain \ELouvainF{} and connected
components \EPlainCCF{}. These rows differ from the proposed method in two
ways at once: how candidates are generated and whether they are scored. With
the same calibrated scorer and an alert threshold selected per method on the
validation seeds, the static rules reach \ECapTenSF{} (cap 10) and
\ECapHundredSF{} (cap 100) and Louvain \ELouvainSF{}, against \EPropF{}
($\pm$\EPropFCI{}) for the proposed method; every difference to the proposed
method remains significant after Holm correction ($p<0.001$). Within the
scored group the differences isolate the candidate generator as a whole
(Section~\ref{sec:attribution} separates its parts); between a generator's
unscored and scored rows they isolate the scorer. The shipped scorer was
trained on the proposed pipeline's candidates, which favours it;
Section~\ref{sec:existing} gives every method its own scorer. For the rules with cap 10, the scorer raises F1
from \ECapTenF{} to \ECapTenSF{}, through precision (\ECapTenP{} to
\ECapTenSP{}) at some cost in recall (\ECapTenR{} to \ECapTenSR{}); the
candidate generator raises it further to \EPropF{}, mostly through recall
(\ECapTenSR{} to \EPropR{}). The scorer was trained on candidates of the proposed pipeline, so
on rule and Louvain candidates it works under some distribution shift. The
unscored static rules reach recall
\ECapTenR{} but precision only \ECapTenP{}: the shared-attribute rules and the
3-cycle rule fire on many benign groups, and two-hit confirmation does not
remove them. Connected components never produce a group within the ring size
limits, because benign hubs join almost all users into one component; Louvain
splits that component, but its communities rarely coincide with a ring. The
thresholds selected on validation seeds improve on the a-priori ones (F1
\EPropF{} vs.\ \EPropDefaultF{}, $p=0.034$, Holm over the eight comparisons
of Table~\ref{tab:baselines}).

Per scenario (Table~\ref{tab:scenarios}), the gap to the static rules is
largest on \texttt{drift} and \texttt{high\_noise}, where many benign groups
satisfy the rules and their precision collapses, and smallest on the clean
\texttt{baseline} scenario and on \texttt{adversarial}, which is the hardest
scenario for every method. Against the scored rules the difference depends on
the scenario: \RfOldNoiseD{} on \texttt{high\_noise}, but only \RfOldHubD{}
(95\% CI \RfOldHubLo{} to \RfOldHubHi{}, wins/ties/losses \RfOldHubWTL{}) on
\texttt{hub\_heavy}, the scenario built for hub hiding. With their own scorer
the rules reach \RfScRulesTabF{} on these streams.

\textbf{Degree cap vs.\ hub weighting.} Fig.~\ref{fig:frontier} isolates the
mechanism this paper replaces. Only shared-attribute and mixed rings count as
ground truth here (flags that match only other ring types are ignored), the
static rules run without the cycle rule, and the hub-aware detector runs alone.
Without the scorer, the rules' recall is \FrCapMaxR{} for every cap from 5 to
200, although caps of 50 and more exceed every benign hub in these scenarios: the
rules miss these rings because each rule returns all users of one attribute,
so a ring inside a larger group fails the Jaccard match, whatever the cap.
Section~\ref{sec:attribution} tests the cap where it does act, inside the
pairwise detector. Hub weighting reaches
recall \FrHubMaxR{}, but on its own at precision \FrHubPAtMaxR{}. With the
calibrated scorer on top, hub weighting reaches recall \FrHubSMaxR{} (95\%
bootstrap interval \RfFrHubSRLo{}--\RfFrHubSRHi{}; precision \FrHubSPAtMaxR{}) and its best F1 is \FrHubSBestF{} (alert threshold
\FrHubSBestFAt{}); the cap with the scorer stays below recall \FrCapSMaxR{}
(\RfFrCapSRLo{}--\RfFrCapSRHi{}) and
its best F1 is \FrCapSBestF{}. The difference in best F1, \FrDiff{}, has a
95\% bootstrap interval of \FrDiffLo{} to \FrDiffHi{} over the 20 runs (the best
threshold re-selected in every resample; whole runs resampled within scenario),
so at the best operating point the two mechanisms are not distinguishable on
F1. The best operating point is chosen on these test runs for both arms, so
both best-F1 values are optimistic. At low recall the cap with the scorer is the
more precise of the two (\FrCapSBestFP{} at recall \FrCapSBestFR{}). Hub
weighting therefore does not dominate the cap everywhere; it extends the
frontier to recall levels the cap cannot reach, and the scorer is what makes
those levels usable.

\begin{figure}[t]
\centering
\includegraphics[width=\columnwidth]{\figpath/fig1_pr_frontier.pdf}
\caption{Precision and recall on shared-attribute rings (micro-averaged over
4 scenarios $\times$ 5 test seeds). Solid: detector alone, swept over the
degree cap or the edge threshold $\theta_e$. Dashed: with the calibrated scorer,
swept over the alert threshold $\theta_a$ (cap 10; $\theta_e=\TunedEdge$).
Best-F1 difference with the scorer (hub minus cap), each at the best alert
threshold on these test runs (an oracle choice): \FrDiff{}, 95\% bootstrap
interval \FrDiffLo{} to \FrDiffHi{} (\texttt{results/e3b\_frontier\_ci}).}
\label{fig:frontier}
\end{figure}

\subsection{Comparison with existing methods}
\label{sec:existing}
% generated by scripts/make_refine_tables.py from results/w2_eval, results/w3b_eval
\begin{table*}[t]
\centering
\caption{Comparison with existing methods on fresh test seeds 2101--2120 (5 scenarios $\times$ 20
seeds, paired). Every method has a scorer trained on its own candidates (seeds 101--110) unless
noted, and its structural parameter and alert threshold selected on seeds 301--303 with at least
the proposed method's budget. PR AUC: ring-level, threshold-free (not defined without a ring
score). $\Delta$: proposed minus row (scenario-weighted, 95\% cluster bootstrap), $p$:
stratified sign-flip; the comparison with the strongest baseline by validation F1 (GraphSAGE, no
scorer) is confirmatory, the others exploratory. $^\ddagger$Exploratory variant added after the
degree-cap control. $^\S$Selection uses the scenario's identity (an oracle advantage).}
\label{tab:existing}
\small
\setlength{\tabcolsep}{2.5pt}
\resizebox{\textwidth}{!}{%
\begin{tabular}{lcccccccccc}
\toprule
Method & P & R & F1 & PR AUC & Hub-hiding & Decoy & Adaptive & $\Delta$ & 95\% CI & $p$ \\
\midrule
Proposed (shipped) & 0.74 & 0.72 & 0.719 $\pm$ 0.040 & 0.780 & 65/105 & 50/59 & 24/28 & -- &  &  \\
Proposed, constant weight$^\ddagger$ & 0.81 & 0.71 & 0.744 $\pm$ 0.043 & 0.755 & 69/105 & 50/59 & 11/28 & $-$0.025 & [$-$0.054, +0.003] & 0.192 \\
\midrule
GraphSAGE, no ring scorer & 0.88 & 0.72 & 0.778 $\pm$ 0.040 & -- & 73/105 & 37/59 & 18/28 & $-$0.059 & [$-$0.097, $-$0.023] & 0.025 \\
GraphSAGE + scorer & 0.85 & 0.69 & 0.751 $\pm$ 0.041 & 0.719 & 69/105 & 37/59 & 16/28 & $-$0.033 & [$-$0.072, +0.007] & 0.227 \\
Rules, cap tuned per scenario$^\S$ & 0.87 & 0.58 & 0.659 $\pm$ 0.053 & 0.619 & 71/105 & 23/59 & 22/28 & +0.059 & [+0.023, +0.094] & 0.014 \\
Rules, cap and cycle rule tuned & 0.91 & 0.55 & 0.646 $\pm$ 0.055 & 0.655 & 54/105 & 23/59 & 18/28 & +0.072 & [+0.037, +0.106] & 0.003 \\
Louvain, resolution tuned & 0.61 & 0.48 & 0.521 $\pm$ 0.051 & 0.562 & 43/105 & 28/59 & 11/28 & +0.198 & [+0.158, +0.239] & $<$0.001 \\
Leiden, resolution tuned & 0.86 & 0.40 & 0.507 $\pm$ 0.054 & 0.678 & 5/105 & 4/59 & 4/28 & +0.212 & [+0.172, +0.252] & $<$0.001 \\
FRAUDAR-style peeling (= Spade output) & 0.85 & 0.37 & 0.456 $\pm$ 0.058 & 0.519 & 58/105 & 2/59 & 5/28 & +0.263 & [+0.221, +0.304] & $<$0.001 \\
Gradient boosting on Louvain groups & 0.91 & 0.33 & 0.427 $\pm$ 0.062 & 0.510 & 0/105 & 7/59 & 3/28 & +0.291 & [+0.249, +0.334] & $<$0.001 \\
\midrule
\multicolumn{11}{l}{\textit{Hybrid, fresh seeds 2301--2320 ($\Delta$: hybrid minus row)}} \\
Proposed + GraphSAGE detector & 0.76 & 0.80 & 0.770 $\pm$ 0.036 & 0.808 & 89/107 & 59/67 & 23/30 & -- &  &  \\
GraphSAGE, no ring scorer & 0.88 & 0.72 & 0.780 $\pm$ 0.037 & -- & 68/107 & 50/67 & 16/30 & $-$0.010 & [$-$0.044, +0.023] & 0.664 \\
GraphSAGE + scorer & 0.84 & 0.70 & 0.755 $\pm$ 0.039 & 0.701 & 69/107 & 47/67 & 14/30 & +0.015 & [$-$0.021, +0.051] & 0.547 \\
Proposed (shipped) & 0.76 & 0.72 & 0.718 $\pm$ 0.037 & 0.792 & 68/107 & 54/67 & 20/30 & +0.052 & [+0.023, +0.081] & 0.003 \\
\bottomrule
\end{tabular}}
\end{table*}

Table~\ref{tab:existing} compares the pipeline with six existing methods under
the protocol of Section~\ref{sec:setup}, on test seeds no experiment had used.
Against every non-learned method it is ahead: by \RfExrulesperscenarioDA{} ring
F1 (95\% CI \RfExrulesperscenarioLo{} to \RfExrulesperscenarioHi{}) against the
rules with a cap and cycle rule tuned per scenario, which is the strongest rule
baseline and uses the scenario's identity, by \RfExrulesDA{} against the rules
tuned globally, and by \RfExlouvainDA{}, \RfExleidenDA{} and \RfExpeelingDA{}
against Louvain, Leiden and FRAUDAR-style peeling. With their own scorers the
rules come much closer than in Table~\ref{tab:baselines} (F1 \RfExrulesF{}
against \ECapTenSF{} with the shipped scorer). Peeling and Leiden are precise
but find few rings, and peeling and Leiden find almost no decoy rings
(\RfExpeelingDecoy{} and \RfExleidenDecoy{}).

A GraphSAGE classifier trained on ring-membership labels is the one method
ahead of the pipeline: F1 \RfExgnnunscoredF{} against \RfExproposedF{}
(difference \RfExgnnunscoredD{}, 95\% CI \RfExgnnunscoredLo{} to
\RfExgnnunscoredHi{}, confirmatory). With a ring scorer of its own GraphSAGE
reaches \RfExgnnF{}, and the difference (\RfExgnnD{}) is not significant. The
gap is not uniform: GraphSAGE is ahead on \texttt{drift} (\RfExgnnunscoredDrift{})
and \texttt{high\_noise} (\RfExgnnunscoredNoise{}), behind on
\texttt{adversarial} (\RfExgnnunscoredAdv{}), and the pipeline finds more
adaptive rings (\RfExproposedAdapt{} against \RfExgnnunscoredAdapt{}) and decoy
rings (\RfExproposedDecoy{} against \RfExgnnunscoredDecoy{}). Ranked by their
calibrated scores, the pipeline's rings have the highest ring-level PR AUC of
all scored methods (\RfExproposedPR{} against \RfExgnnPR{} for GraphSAGE with a
scorer; difference \RtPrD{}, 95\% CI \RtPrLo{} to \RtPrHi{}, exploratory), so
GraphSAGE's advantage lies in its operating point, not in ranking.
All of these data come from one generator, and GraphSAGE learns its regularities
from about 320{,}000 labelled user snapshots.

\textbf{Hybrid.} Adding the GraphSAGE classifier as a fourth detector inside the
pipeline, with a ring scorer trained on the combined candidates, was decided
after this comparison, so we tested it on another fresh block (2301--2320,
bottom of Table~\ref{tab:existing}). The hybrid reaches F1 \RfHyhybridF{}
against \RfHygnnunscoredF{} for GraphSAGE alone (hybrid minus GraphSAGE
\RfHygnnunscoredD{}, 95\% CI \RfHygnnunscoredLo{} to \RfHygnnunscoredHi{}) and
\RfHyproposedF{} for the pipeline alone. Any gain of the hybrid belongs to the
combination, not to a single element of the pipeline.

\subsection{What the advantage depends on}
\label{sec:attribution}
% generated by scripts/make_refine_tables.py from results/r1_cap_control, results/nd_evidence_count
\begin{table*}[t]
\centering
\caption{What the pairwise detector's advantage depends on: isolation controls inside the full
pipeline (same detectors, tracker and scorer; 5 scenarios $\times$ 10 or 20 fresh test seeds).
$\Delta$: scenario-weighted mean paired difference in ring F1 (A minus B) with 95\% cluster
bootstrap interval; $p$: stratified sign-flip test; (H): confirmatory comparison (Holm family of
four); others exploratory. Last column: recall of hub-hiding rings, A / B.}
\label{tab:attribution}
\small
\setlength{\tabcolsep}{3pt}
\resizebox{\textwidth}{!}{%
\begin{tabular}{llcccccc}
\toprule
A vs.\ B & Seeds & F1 A / B & $\Delta$ & 95\% CI & $p$ & $\Delta$ hub\_heavy & Hub-hiding \\
\midrule
Keep hub links vs.\ cap 10, $1/\ln(1+d)$ (H) & 2001--10 & 0.700 / 0.616 & +0.084 & [+0.048, +0.121] & 0.002 & +0.34 & 29/52 / 9/52 \\
Keep hub links vs.\ cap 10, constant & 2001--10 & 0.729 / 0.671 & +0.057 & [+0.039, +0.077] & 0.007 & +0.31 & 31/52 / 7/52 \\
Hub cannot suffice vs.\ can, constant & 2401--20 & 0.737 / 0.600 & +0.137 & [+0.106, +0.168] & $<$0.001 & +0.44 & 72/104 / 9/104 \\
Rare attribute cannot suffice, $1/\ln$ (H) & 2401--20 & 0.729 / 0.729 & 0.000 & [$-$0.028, +0.028] & 0.980 & $-$0.03 & 68/104 / 67/104 \\
$1/\ln(1+d)$ vs.\ constant, no cap & 2001--10 & 0.700 / 0.729 & $-$0.028 & [$-$0.068, +0.009] & 0.218 & 0.00 & 29/52 / 31/52 \\
$1/\ln$ vs.\ constant, both two-evidence & 2401--20 & 0.729 / 0.737 & $-$0.008 & [$-$0.020, +0.003] & 0.174 & $-$0.03 & 68/104 / 72/104 \\
Pairwise + cap 10 vs.\ rules + scorer & 2001--10 & 0.616 / 0.611 & +0.005 & [$-$0.039, +0.051] & 0.885 & $-$0.30 & 9/52 / 27/52 \\
\bottomrule
\end{tabular}}
\end{table*}

Table~\ref{tab:attribution} isolates the parts of the pairwise rule inside the
full pipeline, on fresh seeds. Removing hub links with a degree cap of 10 costs
\RfCapD{} F1 (95\% CI \RfCapLo{} to \RfCapHi{}, confirmatory, $p=\RfCapP$),
\RfCapHub{} on \texttt{hub\_heavy}, where recall of hub-hiding rings falls from
\RfCapHHNo{} to \RfCapHHCap{}. Keeping hub links but letting a single attribute
link a pair (a constant weight with $\theta_e=0.5$) costs \RfNdOneD{}, and
\RfNdOneHub{} on \texttt{hub\_heavy}: every co-user of the hub is linked and the
ring cannot be separated from them. Under the condition that a busy attribute
cannot suffice alone, the weight's shape does not matter: $1/\ln(1+d)$ and a
constant weight differ by \RfLogConstD{} (CI \RfLogConstLo{} to \RfLogConstHi{})
without a cap, and by \RfNdShapeD{} (CI \RfNdShapeLo{} to \RfNdShapeHi{}) when
both also bar rare attributes from linking alone. That last restriction
(the confirmatory test of the evidence condition) changes nothing overall
(\RfNdPrimD{}, CI \RfNdPrimLo{} to \RfNdPrimHi{}) but trades scenarios: it removes
almost all false rings on \texttt{drift} (\RfNdDriftFalseOne{} to
\RfNdDriftFalseTwo{} per run; F1 \RfNdDriftD{}) and loses adaptive rings
(\RfNdAdaptOne{} to \RfNdAdaptTwo{}; \texttt{adversarial} \RfNdAdvD{}), which
often link through one rarely shared device. Finally, with the same cap the
pairwise pipeline and the scored rules tie overall (\RfPairCapVsRulesD{}) but
trade sharply: the pipeline wins under benign noise (\RfPairCapVsRulesNoise{} on
\texttt{high\_noise}) and the rules win on hub hiding
(\RfPairCapVsRulesHub{} on \texttt{hub\_heavy}). The gain over the rules in
Table~\ref{tab:baselines} therefore has two sources: pairwise evidence with
several detectors, which helps under noise, and keeping hub links, which helps
under hub hiding.

Two exploratory follow-ups on further fresh seeds (protocol in
\texttt{docs/research/round3\_protocol.md}) sharpen this. First, the cost of a cap
does not come only from attributes that could link a pair alone. A cap of 27
removes exactly the attributes that, at $\theta_e=0.3$, can never link a pair
without a second piece of evidence, and still costs \RtCapTwentySevenD{} F1 (CI
\RtCapTwentySevenLo{} to \RtCapTwentySevenHi{}; cap 50: \RtCapFiftyD{}; cap 10:
\RtCapTenD{}). Hub-hiding rings found fall from \RtCapHHNo{} to
\RtCapHHTwentySeven{}. Busy attributes therefore help as \emph{supporting}
evidence, added to a transfer or a rarer attribute. Second, the property that
rare attributes link alone is what the adaptive rings need, but a step function
with that property is not a substitute for the discount: $h(d)=1$ for $d\le27$ and
0.2 above costs \RtStepD{} (CI \RtStepLo{} to \RtStepHi{}), mostly on
\texttt{hub\_heavy}, where the flat 0.2 lets busy attributes add up. A constant
weight that requires two pieces of evidence ($\theta_e=1.0$, the constant row of
Table~\ref{tab:existing}) differs by \RtConstTwoD{} (CI \RtConstTwoLo{} to \RtConstTwoHi{}),
with fewer false rings, but finds \RtConstTwoAdapt{} adaptive rings against
\RtLogAdapt{} for $1/\ln(1+d)$.

\subsection{RQ2: latency and scale}
% generated from results/e4_scale/results.csv (sha256 of LF text 7eb8cf45e1ba0a9a, 30 rows)
\begin{table*}[t]
\centering
\caption{Per-batch detection latency (median, ms; mean $\pm$ 95\% CI over 3 seeds)
and users examined per batch, incremental vs full recomputation, in-memory store. With three seeds
the $t$-interval can exceed the mean; the text uses ranges over seeds.}
\label{tab:scale}
\begin{tabular}{rccccc}
\toprule
& \multicolumn{2}{c}{Incremental} & \multicolumn{2}{c}{Full} & \\
Users & Region & ms & Region & ms & Speed-up \\
\midrule
1,000 & 976 & 211.0 $\pm$ 142.9 & 1015 & 159.7 $\pm$ 64.4 & 0.8$\times$ \\
5,000 & 1878 & 177.8 $\pm$ 46.7 & 5016 & 589.8 $\pm$ 67.3 & 3.3$\times$ \\
20,000 & 1324 & 54.6 $\pm$ 37.1 & 20010 & 2683.7 $\pm$ 1062.0 & 49.2$\times$ \\
50,000 & 1110 & 38.5 $\pm$ 38.2 & 50011 & 6181.6 $\pm$ 2576.4 & 160.4$\times$ \\
100,000 & 827 & 36.1 $\pm$ 44.2 & 100013 & 22500.5 $\pm$ 10278.5 & 623.0$\times$ \\
\bottomrule
\end{tabular}
\end{table*}

\begin{figure}[t]
\centering
\includegraphics[width=\columnwidth]{\figpath/fig2_latency_scaling.pdf}
\caption{Median per-batch detection latency as the graph grows (mean over 3
seeds; bars: per-seed minimum and maximum), incremental (dirty region) vs.\ full recomputation of the same
detectors, in-memory store.}
\label{fig:scale}
\end{figure}

Table~\ref{tab:scale} and Fig.~\ref{fig:scale} grow the user population of the
\texttt{scale\_sweep} scenario at a fixed event rate (100 events/s) and compare
incremental detection with recomputing the same detectors on the whole graph
for every batch. Full recomputation grows at least linearly with the graph
(faster between 50{,}000 and 100{,}000 users, where the 7.7\,GB test machine was
short of memory; we did not separate the two effects). The incremental region
stays between a few hundred and two thousand users, because the fan-out limits
stop the expansion at hubs; it even shrinks as the graph grows, since this
scenario scales the number of devices, cards and IPs with the users (so an
ordinary attribute is shared by fewer users) while hub sizes stay fixed. Its
latency therefore stays flat: at \ScaleMaxUsers{} users a batch takes \ScaleIncMs{}\,ms
incrementally against \ScaleFullMs{}\,ms in full. This comparison is a worst-case
bound, not a deployment comparison: nobody would recompute the whole graph on
every 50-event batch, the incremental figure (3 seeds) varies more than its
mean, and the 50{,}000 and 100{,}000-user full figures are confounded by memory
pressure. At
1{,}000 users the region already covers almost the whole graph, and the
incremental mode is slower than full recomputation because of the cost of
expanding the region. On the accuracy scenarios, incremental and full
recomputation flag exactly the same rings in every run (ablation
\emph{no incremental} in Table~\ref{tab:ablation}), so restricting detection to
the dirty region loses nothing there.

\textbf{Against a timer.} The realistic alternative is to run the same
detectors on the whole graph on a timer. On the accuracy scenarios (30 to 200
users, \TmRuns{} paired runs) a 1\,s timer matches the incremental engine in F1
(\TmOneF{} vs.\ \TmIncF{}, $p=\TmOnePval{}$) and time to detect
(\TmOneTTD{}\,s vs.\ \TmIncTTD{}\,s) at lower detection cost (\TmOneCost{}
vs.\ \TmIncCost{}\,ms per simulated second); a 5\,s timer costs less still
(\TmFiveCost{}) but detects later (\TmFiveTTD{}\,s), and a 30\,s timer misses
almost every ring, which ends before two passes confirm it (F1 \TmThirtyF{}).
On small graphs, incremental detection therefore buys nothing. The picture
reverses with size. In \texttt{scale\_sweep} with hubs that grow with the
population (15\% and 8\% of users, other hubs up to 10\%; \GhSeeds{} seeds), the
cost per simulated second is \GhIncCostA{}, \GhIncCostB{} and \GhIncCostC{}\,ms
for incremental detection at 1{,}000, 5{,}000 and 20{,}000 users, against
\GhTimerCostA{}, \GhTimerCostB{} and \GhTimerCostC{}\,ms for a 5\,s timer
(\GhFullCostC{}\,ms for recomputing every batch at 20{,}000). At 5{,}000 users
the two overlap across seeds; at 20{,}000 users the timer needs more than one
second of computation per simulated second in all three seeds, so it cannot
keep up, while the incremental region stays near \GhIncRegionC{} users even
though the hubs grow. On five fresh seeds at 10{,}000 users (revision
container, run alone) the 5\,s timer needs \RfScaleTenTimerLo{}--\RfScaleTenTimerHi{}\,ms
per simulated second, more than one second in \RfScaleTenTimerOver{} of 5 seeds;
the incremental engine is cheaper in \RfScaleTenWins{} of 5 (median
\RfScaleTenIncMed{}\,ms), but in one seed it costs \RfScaleTenIncHi{}\,ms, more
than the timer. The incremental
cost is not monotone in size (\GhIncCostB{} at 5{,}000, \GhIncCostC{} at
20{,}000); a likely cause, which we did not isolate, is that the region is
bounded by the fan-out limits while the number of touched users per batch
falls as the graph grows at a fixed event rate. We did not run 50{,}000 users
with a timer. The static rules run on
the same 5\,s timer, which is why their time to detect is similar to ours on the
accuracy scenarios (\EPropTTD{}\,s vs.\ \ECapTenTTD{}\,s).

\subsection{RQ3: robustness to evasion}
\begin{figure}[t]
\centering
\includegraphics[width=\columnwidth]{\figpath/fig4_hub_size.pdf}
\caption{A shared-IP ring hiding behind an office hub of growing size
(\texttt{hub\_heavy}, 5 test seeds per size). Left: recall of the hub-hiding
rings; right: ring F1 (mean, 95\% CI).}
\label{fig:hubsize}
\end{figure}
\begin{figure*}[t]
\centering
\includegraphics[width=\textwidth]{\figpath/fig5_evasion.pdf}
\caption{Ring recall by evasion strategy, pooled over all runs of
Table~\ref{tab:baselines} (static rules without a scorer; Table~\ref{tab:existing}
gives the recall by strategy for every method with its own scorer).}
\label{fig:evasion}
\end{figure*}

\textbf{Hub size.} In Fig.~\ref{fig:hubsize} the office hub grows from
\HubSizeMin{} to \HubSizeMax{} accounts. Recall of the hub-hiding rings does
not fall with hub size for any method (proposed \HubSizePropHidden{}, static
cap 10 \HubSizeCapTenHidden{} over all sizes). The static rules do not find
these rings through the shared hub: with the 3-cycle rule switched off, cap 10
finds \HubSizeCapTenNoCycleHidden{}, because the rings also move money among
themselves in short cycles. What separates the methods is precision, and so
F1 (\HubSizePropF{} vs.\ \HubSizeCapTenF{} for the unscored rules;
\RfScRulesHubSizeF{} for the rules with their own scorer).

\textbf{Hubs above the fan-out limit.} Fig.~\ref{fig:hubsize} stops at 200
accounts, below the 250-user first-hop limit. We therefore resized the office
hub to 500, 2,000 and 10,000 accounts (5 seeds each) and compared the shipped
engine with full recomputation and with a ``bound'' rule that does cross an
over-limit attribute, towards those co-members whose attribute weight with the
touched user could reach the edge threshold. At every size the bound rule
flagged the same rings as the 250-user limit, with identical precision, recall
and region size (\LhBoundSame{} runs): it never added a user, so no pair in these
streams shared two over-limit attributes. Full recomputation does not give
identical runs. Tracing them showed two causes, neither of them the limit: a
ring none of whose members is touched enters a region only partly, and greedy
peeling of a connected component of several hundred users can return the ring
on the whole graph but not on the region's part of it. In total it found no
more hub-hiding rings (\LhLogFullHHA{} vs.\
\LhLogFanHHA{} at 500, \LhLogFullHHB{} vs.\ \LhLogFanHHB{} at 2,000) and its F1
was not higher (\LhLogFullFA{} vs.\ \LhLogFanFA{}, \LhLogFullFB{} vs.\
\LhLogFanFB{}), while the incremental region stayed at a median of
\LhRegionInc{} users against up to \LhRegionFullMax{}. At 10,000 accounts one
full-recomputation run did not finish within 3.5\,h on the test laptop, so
there is no full-recomputation reference at that size. A constant weight (edge threshold 1.0, as selected in RQ5) found fewer
hub-hiding rings incrementally than with full recomputation at 500 accounts
(\LhConstFanHHA{} vs.\ \LhConstFullHHA{}; F1 \LhConstFanFA{} vs.\
\LhConstFullFA{}, lower in 3 of 5 runs; five seeds cannot reach significance),
by the same two mechanisms; a single busy attribute cannot reach this threshold
on its own. One full run of it took about \LhConstFullWallB{} minutes at 2,000
accounts. So, for hubs up to 10,000 accounts, the 250-user limit itself hid no
pair for either weight, but ring-level equality with full recomputation does
not hold at these sizes. The static rules stay behind (F1 \LhCapTenFA{},
\LhCapTenFB{}, \LhCapTenFC{} without a scorer; \RfScRulesLhA{},
\RfScRulesLhB{}, \RfScRulesLhC{} with their own scorer).

\textbf{Evasion strategies.} Fig.~\ref{fig:evasion} breaks recall down by the
strategy each planted or random ring uses. The proposed method's advantage in
recall comes mostly from rings without evasion (\EvPropNone{} vs.\
\EvCapTenNone{} for cap 10) and rings with decoy edges (\EvPropDecoyEdges{}
vs.\ \EvCapTenDecoyEdges{}). On hub-hiding rings the two are close
(\EvPropHubHiding{} vs.\ \EvCapTenHubHiding{}), and identity rotation
(\EvPropIdentityRotation{}) and low-and-slow activity (\EvPropLowAndSlow{})
defeat every method we tested, presumably because rotation removes the shared
attributes the detectors rely on and slow rings accumulate little evidence
inside one window.

\textbf{Adaptive attacker.} With feedback switched on, the attacker is told
which of its rings were alerted and shifts its strategy mix. Over
\AdaptSeeds{} test seeds it switched strategy \AdaptOnSwitches{} times and
launched \AdaptOnRings{} adaptive rings, of which \AdaptOnFound{} were
detected; without feedback, \AdaptOffFound{} of \AdaptOffRings{}. We saw no
change in the detection rate (difference 0.00, 95\% interval \RfAdaptLo{} to
\RfAdaptHi{}), an interval too wide to rule out an effect, and the attacker's
strategies are limited to the four we implemented.

\subsection{RQ4: sensitivity}
\begin{figure*}[t]
\centering
\includegraphics[width=\textwidth]{\figpath/fig6_sweeps.pdf}
\caption{Ring F1 of the proposed method and of the static rules (cap 10,
without a scorer) when the event rate (\texttt{adversarial}), the planted ring
size (\texttt{baseline}) and the view window (\texttt{drift}, proposed only)
change; 5 test seeds per point. With their own scorer the rules reach F1
\RfScRulesRateLo{} at the lowest and \RfScRulesRateHi{} at the highest event
rate, and \RfScRulesSizeLo{} at the smallest and \RfScRulesSizeHi{} at the
largest ring size.}
\label{fig:sweeps}
\end{figure*}
Raising the event rate from \SwRateLo{} to \SwRateHi{} events/s lowers the
proposed method's F1 from \SwRateLoF{} to \SwRateHiF{} and raises its per-batch
latency from \SwRateLoLat{} to \SwRateHiLat{}\,ms, since each batch touches more
of the graph (Fig.~\ref{fig:sweeps}). Larger planted rings are harder (F1
\SwRingSizeLoF{} at size \SwRingSizeLo{}, \SwRingSizeHiF{} at size
\SwRingSizeHi{}); a likely cause, which we did not isolate, is the 0.75 edge
density the peeling requires. The view window between
\SwWindowLo{} and \SwWindowHi{}\,s changes F1 little (\SwWindowLoF{} vs.\
\SwWindowHiF{}).

\subsection{RQ5: ablations}
% generated from results/e2_ablations/results.csv (sha256 of LF text ac950ecba52299df, 180 rows)
\begin{table*}[t]
\centering
\caption{Ablations on adversarial, hub\_heavy, high\_noise, baseline
($\times$ 5 test seeds). $\Delta$F1: paired mean difference to the full system;
$d_z$: paired effect size of the full system over the variant; $p$: Wilcoxon signed-rank,
uncorrected and Holm-corrected over the 7 tested comparisons.}
\label{tab:ablation}
\begin{tabular}{lcccccccc}
\toprule
Variant & Precision & Recall & F1 & $\Delta$F1 & $d_z$ & $p$ & $p_{\mathrm{Holm}}$ & Latency (ms) \\
\midrule
Full system & 0.75 $\pm$ 0.12 & 0.73 $\pm$ 0.12 & 0.72 $\pm$ 0.11 & -- & -- & -- & -- & 19.2 \\
w/o hub weighting & 0.65 $\pm$ 0.15 & 0.58 $\pm$ 0.14 & 0.60 $\pm$ 0.14 & -0.12 & 0.42 & 0.161 & 0.462 & 24.0 \\
w/o time decay & 0.80 $\pm$ 0.11 & 0.72 $\pm$ 0.10 & 0.75 $\pm$ 0.09 & +0.02 & -0.18 & 0.277 & 0.462 & 17.6 \\
w/o incremental scheduling & 0.75 $\pm$ 0.12 & 0.73 $\pm$ 0.12 & 0.72 $\pm$ 0.11 & 0.00 & -- & -- & -- & 18.4 \\
w/o risk model & 0.28 $\pm$ 0.09 & 0.82 $\pm$ 0.09 & 0.39 $\pm$ 0.10 & -0.33 & 2.05 & $<$0.001 & $<$0.001 & 15.1 \\
w/o hysteresis & 0.72 $\pm$ 0.12 & 0.74 $\pm$ 0.10 & 0.70 $\pm$ 0.10 & -0.02 & 0.37 & 0.154 & 0.462 & 22.7 \\
w/o hub-aware detector & 0.98 $\pm$ 0.05 & 0.25 $\pm$ 0.09 & 0.36 $\pm$ 0.12 & -0.36 & 1.02 & 0.001 & 0.006 & 5.3 \\
w/o cycle detector & 0.75 $\pm$ 0.14 & 0.60 $\pm$ 0.11 & 0.64 $\pm$ 0.11 & -0.09 & 1.02 & 0.002 & 0.009 & 15.4 \\
w/o smurfing detector & 0.72 $\pm$ 0.14 & 0.64 $\pm$ 0.15 & 0.65 $\pm$ 0.14 & -0.07 & 0.93 & 0.003 & 0.013 & 19.9 \\
\bottomrule
\end{tabular}
\end{table*}

Table~\ref{tab:ablation} removes one component at a time (\AblPairs{} paired
runs). The calibrated scorer contributes most: without it F1 falls from
\AblFullF{} to \AblNoScoringF{} because precision drops to \AblNoScoringP{}.
Without the hub-aware detector, F1 falls to \AblNoHubAwareF{} and to
\AblNoHubAwareFHubHeavy{} on \texttt{hub\_heavy}, but precision rises to
\AblNoHubAwareP{} at recall \AblNoHubAwareR{}: most false candidates come from
the hub-aware detector, and the scorer removes most of them. This suggests a
separate alert threshold for rings the hub-aware detector found; selected on
the validation seeds over a $6\times6$ grid, the best pair was the shipped
single threshold (0.8, 0.8), so the refinement gains nothing here. Removing the cycle or the
smurfing detector costs \AblNoCycleF{} and \AblNoSmurfingF{}. Replacing the hub
weight by a constant at the same edge threshold collapses \texttt{hub\_heavy}
(\AblFullFHubHeavy{} to \AblNoHubWeightFHubHeavy{}) but helps
\texttt{high\_noise} (\AblFullFHighNoise{} to \AblNoHubWeightFHighNoise{}). The
effect depends on the scenario (heterogeneity $p=\RfAblHubWHet$), so the pooled
test ($p=\AblNoHubWeightPval{}$) is not informative either way. That ablation keeps the edge threshold tuned for
$1/\ln(1+d)$, which is unfair to the constant; the next paragraph gives each
weight its own threshold.

\textbf{Weight functions.} We compared five weights $h(d)$, each with its edge
threshold selected on the validation seeds and the alert threshold fixed at
\TunedAlert{}: the shipped $1/\ln(1+d)$ ($\theta_e=\WShipEdge{}$), $1/d$
(\WInvEdge{}), $1/\sqrt{d}$ (\WSqrtEdge{}), a soft cap
$1/(1+e^{(d-10)/2})$ (\WSoftEdge{}) and a constant (\WConstEdge{}); for $1/d$ and the soft cap the selected value lies on the edge of this grid
(\RfGridNote{}); for the constant, a wider grid in Section~\ref{sec:attribution}
selects 1.0 again as an interior optimum (1.2 is worse).
On \texttt{hub\_heavy} with \WbSeeds{} paired seeds, the shipped weight reaches
F1 \WbShipF{} and recall \WbShipHH{} of the rings hidden behind hubs; the
constant reaches \WbConstF{} and \WbConstHH{} ($p=\WbConstPval{}$), $1/\sqrt{d}$
\WbSqrtF{}, while $1/d$ (\WbInvF{}) and the soft cap (\WbSoftF{}, hidden-ring
recall \WbSoftHH{} at precision \WbSoftP{}) are worse ($p$ \WbSoftPval{}). Over
the five scenarios the picture is the same: \WcShipF{} for the shipped weight,
\WcConstF{} for the constant and \WcSqrtF{} for $1/\sqrt{d}$, and \WcInvF{} for
$1/d$ and \WcSoftF{} for the soft cap. With the widened grid, $1/d$ at
$\theta_e=0.05$ reaches \RfInvFiveF{} over the five scenarios (\RfInvFiveD{}
against the shipped weight, 95\% CI \RfInvFiveLo{} to \RfInvFiveHi{}) and
\RfInvHubF{} on \texttt{hub\_heavy} (\RfInvHubD{}, CI \RfInvHubLo{} to
\RfInvHubHi{}), so most of its earlier deficit came from the grid edge.
$1/\sqrt{d}$ is equivalent to the shipped
weight within 0.05 F1 (90\% interval \RfWSqrtLo{} to \RfWSqrtHi{}); the constant
is not significantly different but not shown equivalent (\RfWConstLo{} to
\RfWConstHi{}). Switching hub links off loses rings that share a hub
(Section~\ref{sec:attribution}); a strong discount costs little once the edge
threshold is tuned for it, and the soft cap, which nearly switches busy
attributes off, stays behind.
Hysteresis makes no difference beyond 0.05 F1 (90\% interval \RfAblHystLo{} to
\RfAblHystHi{}); time decay makes no significant difference but is not shown
equivalent (\RfAblDecayLo{} to \RfAblDecayHi{}), and without decay F1 is even
slightly higher (\AblNoTimeDecayF{}). The accuracy scenarios last
60--90\,s, short against the 300\,s half-life, so decay has little to act on
here; its value on long streams is untested.

\subsection{Public data}
\label{sec:public}
\input{sections/public.tex}
